\documentclass{article}
\usepackage[T1]{fontenc}
\usepackage{lmodern}
\usepackage{graphicx}
\usepackage{amsmath,amssymb}
\usepackage[a4paper,margin=2cm]{geometry}
\usepackage[absolute,overlay]{textpos}
\setlength{\TPHorizModule}{1mm}
\setlength{\TPVertModule}{1mm}
\usepackage[indonesian]{babel}
\usepackage{hyperref}
\setlength{\emergencystretch}{2em}
% unit: Halaman 25

\begin{document}

% === Halaman 25 (llm) ===

\section*{3. Distinguishing Attack}

\begin{itemize}
    \item Tujuan: menentukan apakah suatu data merupakan keluaran algoritma kriptografi atau data yang benar-benar acak (random).
    \item Jika secara statistik penyerang dapat membedakannya dengan probabilitas yang secara signifikan lebih baik daripada tebakan acak, berarti terdapat \textit{distinguishing attack}.
    \vspace{5mm}
    \item Contoh: misalkan sebuah cipher menghasilkan:
    \[ 10101010101010101010101010... \]
    Pola tersebut jelas tidak menyerupai random.
    Penyerang dapat membuat distinguisher sederhana:
    \textit{Jika jumlah 0 dan 1 sangat tidak seimbang atau terdapat pola tertentu, kemungkinan data tersebut berasal dari cipher.}
    \vspace{5mm}
    \item \textit{Cipher} modern dirancang agar kelemahan seperti ini tidak mudah ditemukan.
\end{itemize}

\end{document}
